\documentclass[10pt,a4paper]{amsart}
\usepackage[margin=19mm]{geometry}
\usepackage{amsmath,amssymb,mathtools}
\usepackage[scr=rsfs,cal=euler]{mathalfa}
\usepackage{xfrac}
\usepackage[unicode,hidelinks,bookmarks=false]{hyperref}
\newcommand{\ie}{i.e.\ }
\newcommand{\cf}{cf.\ }
\newcommand{\resp}{respectively\ }
\newcommand{\Subsection}{\subsection}
\providecommand{\nobreakdash}{\nobreak\hbox{-}}
\setlength{\emergencystretch}{2em}
\multlinegap0pt
\usepackage{dsfont}
\usepackage{float}
\usepackage{tikz}
\usepackage{amssymb}
\newtheorem{definition}{Definition}
\newtheorem{theorem}{Theorem}
\newtheorem{lemma}{Lemma}
\newtheorem{proposition}{Proposition}
\newcommand{\der}{\textup{d}}
\title{Long-time asymptotics for coagulation equations with injection that do not have stationary solutions}
\author{Iulia Cristian, Marina A. Ferreira, Eugenia Franco, Juan J. L. Vel\'{a}zquez}
\date{Source: arXiv:2211.16399v1 (CC BY 4.0)}
\begin{document}
\maketitle

\noindent\emph{Source and licence.} Long-time asymptotics for coagulation equations with injection that do not have stationary solutions. Version arXiv:2211.16399v1 (CC BY 4.0), distributed by arXiv under the Creative Commons Attribution 4.0 International licence, http://creativecommons.org/licenses/by/4.0/, as recorded in the arXiv licence metadata for this exact version and shown on the abstract page https://arxiv.org/abs/2211.16399v1. Full licence notice: \texttt{licenses/arxiv-2211.16399-CC-BY-4.0.txt}.\par\smallskip

\noindent\textit{Editorial scope.} Counted unit: the article's proposition region (source0.tex lines 1456--1467). The article's complete original proof of it is delivered with it (source0.tex lines 1468--1667, one \texttt{proof} environment of 200 lines). No mathematics is written, repaired or re-derived: the statement and the proof are the article's own lines, in the article's own notation, and no retained line is substituted or re-flowed.  Retained uncounted as the source-local context the proof needs: lines 233--237; lines 245--247; lines 250--252; lines 961--963; lines 1010--1013; lines 1021--1026; lines 1034--1048; lines 1055--1060; lines 1126--1146; lines 1133--1136; lines 1383--1385.  Nothing was omitted from any retained span, so the package carries no omission record.  No retained line contains a citation, so the wrapper carries no bibliography and no bibliography entry.  No retained line contains a figure environment, an image-inclusion command or drawing code, so no image is delivered and the package declares no asset dependency.  Editorial formatting only, in addition: an AMS \texttt{amsart} preamble that restates the article's own declarations and macros used by the retained lines, namely the environments \texttt{definition}, \texttt{theorem}, \texttt{lemma}, \texttt{proposition} on separate counters, together with the standard packages those lines need.  The retained environments print in this document as Definition 1, Theorem 1, Lemma 1, Proposition 1 in order of appearance, whereas the article prints the same items with the numbering of its own full document.  The complete native source of the article as distributed in the arXiv e-print for this exact version is bundled unchanged as \texttt{sources/134201/source0.tex}.
\begin{equation} \label{kernel bounds} 
c_2 \left[  \frac{x^{\gamma+\lambda}}{y^{\lambda}}+\frac
{y^{\gamma+\lambda}}{x^{\lambda}} \right]  \leq K\left(  x,y\right)   \leq c_1 \left[  \frac{x^{\gamma+\lambda}}{y^{\lambda}}+\frac
{y^{\gamma+\lambda}}{x^{\lambda}} \right] \textup{ and } \gamma+2\lambda\geq 0,
\end{equation}

\begin{equation}\label{kernel cont}
    K\in\textup{C}((0,\infty)^{2}).
\end{equation}

\begin{equation}\label{avoid_gelation_parameters}
    \gamma<1 \textup{ and } \gamma+\lambda<1.
\end{equation}

\begin{equation}\label{delta1formula}
\rho(C_1):= \left(\frac{1-\gamma }{2 C_1} \right)^{\frac{1}{1-\gamma -\lambda}}.
\end{equation} 

\begin{align}
\label{truncated time dependent eq}
&\partial_{t} \Phi \left( t, \xi \right)- \frac{3+\gamma}{1-\gamma}\Phi(t, \xi) - \frac{2\xi}{1-\gamma} \partial_\xi \Phi(t, \xi) + \frac{\partial_\xi \left( \xi^{\gamma+\lambda} \Phi(t, \xi)\right)  }{\int_{\mathbb R_*} x^{\gamma+\lambda} \Phi(t, \der x ) } =\mathbb K_R [\Phi](t, \xi).
\end{align}

 \begin{equation}\label{truncated kernel} 
 \begin{cases}
 K_R(x,y)=0, &\text{ if } x>R \text{ or } y > R; \\
  | K_R(x,y) - K(x,y) |\leq e^{-R}, &\text{ if } (x,y) \in \left[ \frac{1}{2 R}, \frac{R}{4}\right]^2.
 \end{cases} 
 \end{equation} 

\begin{definition}\label{def:truncated time dep} 
Let $K_R$ be the truncated kernel defined as in \eqref{truncated kernel} as a function of the homogeneous symmetric coagulation kernel $K$ satisfying
\eqref{kernel bounds}, \eqref{kernel cont},  with homogeneity $\gamma < 1 $ and with $\gamma + \lambda < 1$ and $\gamma + 2 \lambda \geq 1 $. A function $\Phi \in C^1([0,T]; \mathcal M_{+,b}(\mathbb R_*))$ is a solution of equation \eqref{truncated time dependent eq} if 
\[
0 < \inf_{  t \in [0, T]  }\int_{\mathbb R_*} x^{\gamma+\lambda } \Phi(t, \der x ) \ \text{ and } \ \sup_{  t \in [0, T]  }\int_{\mathbb R_*} x^{\gamma+\lambda } \Phi(t, \der x ) < \infty 
\] 
and if $\Phi$ satisfies 
\begin{align}\label{weak truncated time dependent eq}
& \int_{\mathbb R_*} \varphi(x) \dot{\Phi}(t, \der x)- \frac{1+\gamma}{1-\gamma} \int_{\mathbb R_*}  \varphi(x) \Phi(t, \der x) +
\frac{2}{1-\gamma} \int_{\mathbb R_*}  \varphi'(x) x \Phi(t, \der x) \\
& - \frac{1}{\int_{\mathbb R_*}x^{\gamma +\lambda}\Phi(t,\der x)}  \int_{\mathbb R_*}  \varphi'(x) x^{\gamma + \lambda}\Phi(t, \der x) \nonumber \\
& =\frac{1}{2}\int_{\mathbb R_*} \int_{\mathbb R_*} K_R(x,y) \left[ \varphi(x+y)- \varphi(x) - \varphi(y) \right] \Phi(t, \der x) \Phi(t, \der y), \nonumber 
\end{align} 
for every test function $\varphi \in C^1_c(\mathbb R_*).$
\end{definition}

\begin{theorem}\label{thm:time dep truncated}
Assume $K_{R}$ to be a truncated kernel defined as in \eqref{truncated kernel} as a function of a homogeneous symmetric kernel $K$ satisfying \eqref{kernel bounds},  \eqref{kernel cont},  with parameters $\gamma, \lambda \in \mathbb R$ such that \eqref{avoid_gelation_parameters} holds and $\gamma + 2\lambda \geq 1 $ and $\gamma > -1.$ 
Then there exist two constants $k_1, k_2>0$ and a constant $R_0>0$ such that if $\Phi_{0} \in S\left(2 k_1 , \frac{k_2}{2} , \rho(k_2)\right) \cap \left\{ H \in \mathcal M_+(\mathbb R_*): H((4R, \infty) )=0 \right\} $ for $R >R_0$ and for $\rho(k_2)$ given by \eqref{delta1formula},
then there exists a unique solution $\Phi \in C^{1}([0,T];\mathcal M_{+}(\mathbb R_*))$ of equation \eqref{truncated time dependent eq} in the sense of Definition \ref{def:truncated time dep}, for $T>0$ sufficiently small.  
For every $ t \in [0, T]$, we have that $\Phi(t, \cdot ) \in S\left(k_1, k_2, \rho( k_2 )\right) \cap \left\{ H \in \mathcal M_+(\mathbb R_*): H((4R, \infty) )=0 \right\}.$ 
\end{theorem}

\begin{lemma}\label{lem:auxiliary time dep truncated} 
Let $K_R$ be a truncated kernel defined as in \eqref{truncated kernel} as a function of a homogeneous symmetric kernel $K$ satisfying \eqref{kernel bounds},  \eqref{kernel cont},  with $\gamma , \lambda \in \mathbb R$ satisfying \eqref{avoid_gelation_parameters}. 
Let $k_1, k_2, R$ be three positive constants. 
Assume that the initial condition $\Phi_0$ is such that $ 2k_1 \leq \int_{\mathbb R_*} x^{\gamma+\lambda} \Phi_0(\der x ) \leq \frac{k_2}{2}$ and that $\Phi_0((0,\rho(k_2)) \cup (4R, \infty ))=0$ for $\rho(k_2)$ given by \eqref{delta1formula}. 

For a sufficiently small time  $T>0$, depending on $R$, 
there exists a function $F \in C([0, T] , \mathcal M_{+}(\mathbb R_*))$, with  $F (0, \cdot)=\Phi_0$, that satisfies  
\begin{align} \label{eq for F}
&\int_{\mathbb R_*} \varphi(x) F(t,  \der x) =  \int_{\mathbb R_*} \varphi(x) \Phi_0( \der x)  \\
&+ \frac{1}{2} \int_0^t \int_{\mathbb R_*} \int_{\mathbb R_*} K_R(X(s,x, \alpha),X(s,y,\alpha))e^{\frac{1+\gamma}{1-\gamma}s}\chi^\alpha_{\varphi}(x,y,s)F(s,  \der x) F(s,  \der y) \der s, \nonumber
\end{align}
 for every $\varphi \in C_c(\mathbb R_*)$, with
 \begin{equation}\label{alfa}
\alpha(t):=e^{\frac{1+\gamma}{1-\gamma} t}\int_{\mathbb R_*} \left(X(t,x, \alpha)\right)^{\gamma+\lambda} F(t, \der x ), \quad \forall t \in [0,T].
 \end{equation} 

The function $F$ is such that $ k_1 \leq \alpha (t) \leq k_2$ and such that
\begin{equation} \label{support FPhi}
F(t,(0,\rho(k_2) )\cup (4R, \infty ))=0, \quad  \forall t \in[0,T].
\end{equation} 
\end{lemma}

\begin{align} 
&\int_{\mathbb R_*} \varphi(x) F(t,  \der x) =  \int_{\mathbb R_*} \varphi(x) \Phi_0( \der x)  \\
&+ \frac{1}{2} \int_0^t \int_{\mathbb R_*} \int_{\mathbb R_*} K_R(X(s,x, \alpha),X(s,y,\alpha))e^{\frac{1+\gamma}{1-\gamma}s}\chi^\alpha_{\varphi}(x,y,s)F(s,  \der x) F(s,  \der y) \der s, \nonumber
\end{align}

\begin{align}\label{phi function of F}
\int_{\mathbb R_*} \varphi(x) \Phi(t, \der x)=\int_{\mathbb R_*} \varphi(X(t,y, \alpha))  e^{\frac{1+\gamma}{1-\gamma }t} F(t, \der y),
\end{align} 

\begin{proposition}[Invariant region] \label{prop:inv region}
Assume $K_R $ to be a truncated kernel as in \eqref{truncated kernel} defined as a function of the homogeneous symmetric kernel $K$ that satisfies \eqref{kernel bounds},  \eqref{kernel cont},  for parameters $\gamma, \lambda \in \mathbb R$ such that \eqref{avoid_gelation_parameters} holds and such that $\gamma + 2\lambda \geq  1$, $\gamma >-1 $. Then there exist some positive constants $C_1, C_2$ that do not depend on the truncation $R$ such that the set
\begin{align} \label{invariant region}
P:=\left\{ 
H \in \mathcal M_{+}(\mathbb R_*) : \quad \begin{aligned}   & M_1(H )=1,\ 
H((0, \rho(C_1)) \cup (4R,\infty) )=0, \\
& \frac{1}{C_{2}} \leq M_{\gamma + \lambda}(H) \leq C_1
 \end{aligned} 
 \right\}
\end{align} 
is invariant under the evolution operator $S(t),$ where $\rho(C_1)$ was defined in \eqref{delta1formula}. 
\end{proposition}

\begin{proof}

Let $\Phi_0$ and $K_R$ be as in Lemma \ref{lem:auxiliary time dep truncated}. 
Let $T$ be as in Lemma \ref{lem:auxiliary time dep truncated}. We prove that there exist two constants $C_{1},C_{2}>0$ that do not depend on the truncation parameter $R$ such that the solution $\Phi$ obtained in Theorem \ref{thm:time dep truncated} is such that for every time $t \in [0, T] $ 
\begin{itemize}
        \item[(1)]
$\int_{\mathbb R_* } x \Phi(t,\der x) = 1$; 
    \item[(2)]$\int_{\mathbb R_* } x^{\gamma + \lambda } \Phi(t, \der x) \leq \max\{ C_1,\int_{\mathbb R_* } x^{\gamma + \lambda } \Phi_{0}( \der x)\}$; 
    \item[(3)] $\int_{\mathbb R_* } x^{2- \gamma - \lambda } \Phi(t, \der x) \leq  \max\{C_{2},\int_{\mathbb R_* } x^{2-\gamma-\lambda} \Phi_{0}( \der x)\}$.
\end{itemize}
Additionally, we prove that we can conclude from (2) that $\Phi(t,(0,\rho(C_{1}))=0$, for every $t\in[0,T]$.

We first notice that, since $\Phi(t, \cdot)$ has compact support, we can consider in equation \eqref{weak truncated time dependent eq} test functions $\varphi$ such that $\varphi(x)=x^{k} $, for $k\in\mathbb{R}.$

Let us prove (1). We consider in equation \eqref{weak truncated time dependent eq} a test function $\varphi$ such that $\varphi(x)=x $.
We deduce that 
\[
\frac{d}{dt} \int_0^\infty \xi \Phi(t, \der \xi)=1- \int_0^\infty \xi \Phi(t,\der \xi).
\]
Hence 
\[
\int_0^\infty \xi \Phi(t,\xi) d \xi = 1+ \left( \int_{\mathbb R_*} \xi \Phi_0(\der\xi) -1 \right) e^{-t} 
\]
which since $\int_{\mathbb R_*} \xi \Phi_0(\der \xi)=1 $ implies 
\begin{equation} \label{lower bound mass truncated}
 \int_{\mathbb R_*}\xi \Phi(t, \der \xi) = 1.
\end{equation} 



We prove (2). 
We start by proving this for $\gamma + 2 \lambda >1. $
Consider a test function $\varphi $ such that $\varphi(x)=x^{\gamma +\lambda}$
in equation \eqref{weak truncated time dependent eq}. 
Then, we can see that the moment $M_{\gamma + \lambda}$ satisfies the following 
\begin{align*}
& \partial_{t}M_{\gamma +\lambda} \leq - \frac{2\lambda + \gamma -1 }{1-\gamma} M_{\gamma+\lambda}
+ \frac{(\gamma +\lambda)}{M_{\gamma+\lambda}}  M_{2(\gamma+\lambda)-1} \\
&+ c_{3}\int_{[\rho(2C_{1}), \frac{R}{4}] }
 \int_{[\rho(2C_{1}), \frac{R}{4} ] } (x^{\gamma + \lambda} y^{-\lambda} + x^{-\lambda} y^{\gamma + \lambda} ) [ (x+y)^{\gamma + \lambda} - x^{\gamma + \lambda} - y^{\gamma + \lambda}]  \Phi(t, \der x ) \Phi(t, \der y ) 
 \end{align*} 
since, due to the definition of $K_{R}$ in (\ref{truncated kernel}), we have that on the set $[\rho(2C_{1}),\frac{R}{4}]^{2}$, there exists a constant $C>0$ such that $K_{R}\geq\frac{K}{C}$. 

We denote by $z:=\frac{y}{x}$. Notice that since $\gamma + \lambda >0$ when $y \leq x $ we have that 
\begin{align*}
 \left[x^{\gamma + \lambda} y^{-\lambda} + x^{-\lambda} y^{\gamma + \lambda} \right] \left[ (x+y)^{\gamma + \lambda} - x^{\gamma + \lambda} - y^{\gamma + \lambda} \right]& \leq 2  x^{2(\gamma + \lambda)} y^{-\lambda} [ (1+z)^{\gamma + \lambda} - 1- z^{\gamma + \lambda}  ]\\
& \leq 2 x^{2(\gamma + \lambda)} y^{-\lambda} \left( \left(\gamma + \lambda \right)z- z^{\gamma + \lambda} \right)\\
&\leq 2 (\gamma + \lambda -1 ) y^{\gamma} x^{\gamma + \lambda}. 
\end{align*}
As a consequence, by symmetry, we deduce that 
 \begin{align} \label{epsilon appears}
&\partial_{t}M_{\gamma +\lambda}  \lesssim - \frac{2\lambda + \gamma -1 }{1-\gamma} M_{\gamma+\lambda}
+ \frac{(\gamma +\lambda)}{M_{\gamma+\lambda}}  M_{2(\gamma+\lambda)-1} - 2c_{3} \left(1-\gamma-\lambda \right)
 M_{\gamma } M_{\gamma+ \lambda}\nonumber\\
 &+2c_{3}(1-\gamma-\lambda)\int_{(\frac{R}{4},\infty)}x^{\gamma}\Phi(\der x)M_{\gamma+\lambda}+2c_{3}(1-\gamma-\lambda)\int_{(\frac{R}{4},\infty)}x^{\gamma+\lambda}\Phi(\der x)M_{\gamma}\nonumber\\
 &\leq - \frac{2\lambda + \gamma -1 }{1-\gamma} M_{\gamma+\lambda}
+ \frac{(\gamma +\lambda)}{M_{\gamma+\lambda}}  M_{2(\gamma+\lambda)-1} -2c_{3} \left(1-\gamma-\lambda \right)
 M_{\gamma } M_{\gamma+ \lambda}\nonumber\\
 &+4^{1-\gamma}R^{\gamma-1}2c_{3}(1-\gamma-\lambda)M_{\gamma+\lambda}+c_{4}M_{\gamma}\nonumber\\
 &= - \big[\frac{2\lambda + \gamma -1 }{1-\gamma}-\tilde{\varepsilon}\big]M_{\gamma+\lambda}
+ \frac{1}{M_{\gamma+\lambda}} [ (\gamma +\lambda) M_{2(\gamma+\lambda)-1} - 2c_{3} \left(1-\gamma-\lambda \right)
 M_{\gamma } M_{\gamma+ \lambda}^2\nonumber\\
 &+c_{4}M_{\gamma+\lambda}M_{\gamma}],
\end{align}
with $\tilde{\varepsilon}:=4^{1-\gamma}R^{\gamma-1}2c_{3}(1-\gamma-\lambda)$. Notice we can take $\tilde{R}$ sufficiently large so that $- \frac{2\lambda + \gamma -1 }{1-\gamma}+\tilde{\varepsilon}<0$, for every $R\geq \tilde{R}$.

Since $\gamma < 2(\gamma +\lambda ) -1 \leq 1 $ we deduce, by interpolation, using that $M_{1}(\Phi(t))=1$, for all $t\in[0,T]$, that there exists two positive constants $c_5$ and $c_6$ such that 
\[
M_{ 2(\gamma +\lambda ) -1 } \leq c_5 M_\gamma  + c_6. 
\]
Hence, since $\gamma + \lambda >0$ then 
\begin{align*}
\partial_{t} M_{\gamma +\lambda}& \leq- \big[\frac{2\lambda + \gamma -1 }{1-\gamma}-\tilde{\varepsilon}\big] M_{\gamma+\lambda}
+ \frac{1}{M_{\gamma+\lambda}} [ (\gamma +\lambda) \left[ c_5 M_\gamma  + c_6 \right]\\
&- 2c_3 \left(1-\gamma-\lambda \right)
 M_{\gamma } M_{\gamma+ \lambda}^2 +c_{4}M_{\gamma+\lambda}M_{\gamma}].
\end{align*}
Multiplying by $M_{\gamma +\lambda}$ the inequality we deduce that 
\begin{align*}
  M_{\gamma +\lambda}  \partial_{t} M_{\gamma +\lambda} \leq&- \big[\frac{2\lambda + \gamma -1 }{1-\gamma}-\tilde{\varepsilon}\big] M_{\gamma+\lambda}^2
+  [ (\gamma +\lambda) \left[ c_5 M_\gamma  + c_6 \right]\\
&- 2c_3 \left(1-\gamma-\lambda \right)
 M_{\gamma } M_{\gamma+ \lambda}^2+c_{4}M_{\gamma+\lambda}M_{\gamma}]
\end{align*}
which readjusting the constants implies 
\begin{align}\label{constant gamma lambda larger one}
&  \partial_{t} M_{\gamma +\lambda}^2 \leq- c_3 M_{\gamma+\lambda}^2
+  M_\gamma \left( c_4 +c_{5}M_{\gamma+\lambda}  - c_6 M_{\gamma+ \lambda}^2 \right) + c_7
\end{align}
for $c_3, c_4, c_5, c_6,c_{7}>0.$
This implies that the set $\left\{ \Phi \in \mathcal M_{+}(\mathbb R_*) : M_{\gamma +\lambda } \leq  \frac{ c_5+\sqrt{c_5^{2}+4c_{4}c_{6}}}{2 c_{6} } \right\} $ is invariant when $\gamma + 2 \lambda >1.$

We consider now the case $\gamma + 2 \lambda =1.$ 
First of all notice that for every $M < R $ we have 
\[
\int_{[M, \infty ) } x^{\gamma + \lambda} \Phi(t, \der x) \leq M^{\gamma+\lambda -1 } \int_{[M , \infty )} x \Phi(t, \der x ) \leq M^{\gamma + \lambda -1}.
\] 

Notice that an upper bound is obvious for the points $t\in[0,T]$ for which $M_{\gamma+\lambda}(\Phi(t))\leq 1$. 
If there exist $t_{1},t_{2}\in[0,T]$, $t_{1}<t_{2}$, such that  $M_{\gamma+\lambda}(\Phi(t_{1}))<1$ and $M_{\gamma+\lambda}(\Phi(t_{2}))>1$, by the continuity in time of $\Phi,$ we have that there exists $\overline{t}\in[t_{1},t_{2}]$ such that $M_{\gamma+\lambda}(\Phi(\overline{t}))=1$ and $\varepsilon_{1}\in(0,1)$ such that $M_{\gamma+\lambda}(\Phi(s))\geq 1$ on $[\overline{t},\overline{t}+\varepsilon_{1}]$.

On the interval $[\overline{t},\overline{t}+\varepsilon_{1}]$, we can apply the following logic.

We select $R$ large enough so that we can select $M$ to be such that $(1-\delta)^{\frac{1}{1-\gamma-\lambda}} < M < R$, but independent on $R$, and we deduce that there exists a $\delta >0$ such that
\begin{align*}
\int_{(0, M ) } x^{\gamma + \lambda}  \Phi(t,\der x)
&= \int_{\mathbb R_* } x^{\gamma + \lambda}\Phi(t, \der x) - \int_{[M, \infty ) } x^{\gamma + \lambda} \Phi(t, \der x) \\
&\geq \int_{\mathbb R_* } x^{\gamma + \lambda}\Phi(t, \der x)- M^{\gamma + \lambda -1 }\geq 1- M^{\gamma + \lambda -1 } \geq \delta. 
\end{align*}
As a consequence we deduce that 
\begin{equation} \label{gamma bound}
M_\gamma  \geq \int_{(0, M)} x^{\gamma} \Phi(t, \der x )  \geq M^{-\lambda} \int_{(0,M)} x^{\gamma+\lambda} \Phi(t, \der x)  \geq  \delta M^{- \lambda}. 
\end{equation}
Substituting the test function $\varphi(x)=x^{\gamma + \lambda} $ in equation \eqref{weak truncated time dependent eq},
we deduce that there exists a constant $c>0$ such that
\begin{align*}
&\partial_t M_{\gamma + \lambda } \leq \frac{\gamma + \lambda}{M_{\gamma+\lambda} } M_{2(\gamma + \lambda ) -1 } - c (1-\gamma -\lambda ) M_{\gamma +\lambda} M_\gamma+ \tilde{\epsilon}  M_{\gamma +\lambda} + c M_\gamma, 
\end{align*}
for $\tilde{\varepsilon}\in(0,1)$, which can be made sufficiently small as before. Hence, since for suitable constants $c_3, c_4 >0$ we have that $M_{2(\gamma+\lambda )-1} \leq c_3 M_\gamma + c_4$, similarly to (\ref{epsilon appears}), we deduce that
\begin{align*}
&\frac{1}{2}\partial_t M^2_{\gamma + \lambda } 
\leq c_5 M_{\gamma }   (\gamma + \lambda)  +(\gamma + \lambda) c_6 - c_7 M^2_{\gamma + \lambda } M_\gamma+\tilde{\varepsilon} M_{\gamma +\lambda}^{2}+c_{8} M_\gamma M_{\gamma+\lambda}.
\end{align*}
 Using \eqref{gamma bound} we deduce that there exists a constant $c_{9}>0$ such that
\begin{align*}
    &\frac{1}{2}\partial_t M^2_{\gamma + \lambda } 
\leq M_{\gamma }  \left[c_5   (\gamma + \lambda)  +(\gamma + \lambda) c_{9}  +c_{8}M_{\gamma+\lambda}-(c_7-\tilde{\varepsilon}\delta^{-1} M^{\lambda}) M^2_{\gamma + \lambda }  \right] .
\end{align*}
Choosing now $\tilde{\varepsilon}$ sufficiently small such that $-c_7+\tilde{\varepsilon}\delta^{-1} M^{\lambda}<0$, we have that there exists a constant $C_{1}>0$ such that 
\begin{align}\label{constant for gamma lambda}
\int_{\mathbb R_*} x^{\gamma + \lambda } \Phi(t, \der x ) \leq C_{1}. 
\end{align}

We now prove that $\left(S(t) \Phi_0 \right) ((0, \rho(C_1)))=0$. 
Notice that, by Lemma \ref{lem:auxiliary time dep truncated}, we know that $\left(S(t) \Phi_0 \right) ((0, \rho(2C_{1})))=0$, where we recall that $\rho( 2 C_1) \leq \rho(C_1)$. 
Consider now a sequence of test functions $ \{\varphi_n \}$ such that, for every $n\in\mathbb{N}$, $\varphi_{n}$ is decreasing and supported in $(0, \rho(C_1))$. Since $\varphi_n $ is decreasing, we deduce that the solution $F$ of equation \eqref{eq for F} satisfies 
\begin{align*} 
 \int_{\mathbb R_*} \varphi_n(x) F(t, \der x ) \leq \int_{\mathbb R_*} \varphi_n(x) \Phi_{0}(\der x )+ C(R)  \int_{0}^{t} \int_{\mathbb R_*} \varphi_n(x) F(s, \der x ) \der s, \quad \text{ for every } n \in \mathbb N, 
\end{align*} 
where we used in addition that $M_{1}(F)$ is bounded. Using Gr\"onwall’s inequality, we deduce that
\begin{align*}
    \int_{\mathbb R_*} \varphi_n(x) F(t, \der x ) =0,
\end{align*}
for every $t\in[0,T]$ and $n\in\mathbb{N}$.


We then let $\varphi_{n}$ converge pontwise to $\chi_{(0, \rho(C_1))}$ and obtain that
\begin{align*}
F(t, (0, \rho(C_1))) =0,
\end{align*}
for every $t\in[0,T]$.

Consider now a test function $\varphi$ such that $\varphi(x)=0$ for every $x \geq \rho(C_1)$ in equality \eqref{phi function of F}. Then 
\[
\int_{\mathbb R_* }\varphi(x) \Phi(t, \der x )= e^{\frac{1+\gamma}{1-\gamma}t} \int_{[\rho(C_1), \infty) }\varphi(X(t, y, \alpha) ) F(t, \der y). 
\] 
Using the fact that $X(t, y, \alpha) \geq \rho(C_1)$ for every $y \geq \rho(C_1)$, we deduce that 
$\Phi(t, (0, \rho(C_1)))=0$, for every $t \in [0, T]$. 

We conclude by proving (3). 
We consider a test function $\varphi $ equal to $x^{2-\gamma - \lambda} $ in equation \eqref{weak truncated time dependent eq} and deduce that
\begin{align} \label{eq weak form x 2-lambda-gamma} 
&\partial_t \int_{\mathbb{R}_{\ast}} \xi^{2-\gamma - \lambda }\Phi(t,\der \xi) - \frac{3-3\gamma- 2 \lambda}{1-\gamma}\int_{\mathbb{R}_{\ast}}\xi^{2-\gamma - \lambda }\Phi(t,\der \xi) 
- \frac{ 2-\gamma -\lambda }{M_{\gamma+\lambda}(\Phi(t)) } M_{1}(\Phi(t))  \\
&= \frac{1}{2}\int_{[ \rho(C_1), \infty)}\int_{[ \rho(C_1), \infty)} K_R(\xi,z) \left[(z+\xi)^{2-\gamma -\lambda }- \xi^{2-\gamma -\lambda}-z^{2-\gamma -\lambda} \right] \Phi(t,\der \xi)\Phi(t,\der z).\nonumber 
\end{align} 
By Cauchy-Schwarz inequality it follows that
\begin{equation}\label{eq following from CS}
\frac{1}{\int_{[ \rho(C_1), \infty)} \xi^{\gamma +\lambda} \Phi(t,\der \xi)} \leq 
\frac{\int_{[ \rho(C_1), \infty)}\xi^{2-\gamma - \lambda } \Phi(t,\der \xi)}{ \left( \int_{[ \rho(C_1), \infty)}\xi \Phi(t,\der \xi) \right)^2} = \int_{[ \rho(C_1), \infty)} \xi^{2-\gamma - \lambda } \Phi(t,\der \xi). 
\end{equation}
Plugging (\ref{eq following from CS}) into equation \eqref{eq weak form x 2-lambda-gamma} and using the fact that the total mass is equal to $1$ proved in \eqref{lower bound mass truncated}, we deduce that
\begin{align*}
&\partial_t \int_{[ \rho(C_1), \infty)} \xi^{2-\gamma - \lambda }\Phi(t,\der \xi) \leq \left( \frac{3\gamma-3+2\lambda}{1-\gamma} + 2-\gamma -\lambda\right) \int_{[ \rho(C_1), \infty)} \xi^{2-\gamma - \lambda }\Phi(t,\der \xi) \\
&+ \frac{1}{2}\int_{[ \rho(C_1), \infty)} \int_{[ \rho(C_1), \infty)}K_R(\xi,z) \left[(z+\xi)^{2-\gamma -\lambda }- \xi^{2-\gamma -\lambda}-z^{2-\gamma -\lambda} \right] \Phi(t,\der \xi)\Phi(t,\der z).
\end{align*} 

Hence, since $-1 < \gamma <1$ and $0 < \gamma + \lambda <1 $, we have
\begin{align*}
    \frac{3\gamma-3+2\lambda}{1-\gamma} + 2-\gamma -\lambda=- \frac{(\gamma+1)(1-\gamma-\lambda)}{1-\gamma } <0.
    \end{align*}
By symmetry, 
\begin{align*}
   &\frac{1}{2}\int_{[ \rho(C_1), \infty)} \int_{[ \rho(C_1), \infty)}K_R(\xi,z) \left[(z+\xi)^{2-\gamma -\lambda }- \xi^{2-\gamma -\lambda}-z^{2-\gamma -\lambda} \right] \Phi(t, \der\xi)\Phi(t,\der z) \\
   &\leq  \int_{[ \rho(C_1), \infty)} \int_{[\rho(C_1), z]} K_R(\xi,z) \left[(z+\xi)^{2-\gamma -\lambda }- \xi^{2-\gamma -\lambda}-z^{2-\gamma -\lambda} \right] \Phi(t,\der  \xi)\Phi(t,\der z).
\end{align*}
Assume $\xi\leq z$.  Denote $\eta:=\frac{\xi}{z}\in(0,1]$ and observe
\begin{align}\label{moments larger than one}
K_{R}(\xi,z)[(\xi+z)^{2-\gamma-\lambda}-\xi^{2-\gamma-\lambda}-z^{2-\gamma-\lambda}]\leq & K(\xi,z) z^{2-\gamma-\lambda}[(1+\eta)^{2-\gamma-\lambda}-1-\eta^{2-\gamma-\lambda}]\nonumber\\
\leq & C   K(\xi,z)z^{2-\gamma-\lambda}\eta \nonumber \nonumber\\
\leq& C z^{\gamma}(\eta^{\gamma+\lambda}+\eta^{-\lambda}) z^{2-\gamma-\lambda}\eta \leq 2C z^{2-\lambda}\eta^{1-\lambda}\nonumber \\
\leq & 4C (z\xi^{1-\lambda}+\xi z^{1-\lambda}).
\end{align}

Since $\rho(C_1)\leq \xi$ and $\rho(C_1)\leq z$, then $z^{1-\lambda}\leq\rho(C_1)^{-\lambda}z$. Hence 
\[
\frac{d }{dt } \int_{\mathbb R_*} \xi^{2-\gamma-\lambda} \Phi(t, \der \xi) \leq - c_3 \int_{\mathbb R_*} \xi^{2-\gamma-\lambda} \Phi(t, \der \xi) + c(\rho(C_1)),
\]
for suitable constants $c_3, c(\rho(C_1))>0$. Then, (3) follows.
\end{proof}

\end{document}
